\documentclass[12pt]{report}
\usepackage{comment}
\usepackage{xspace}
\usepackage{enumitem}
\usepackage{hyperref}
\usepackage{listings}
\usepackage{graphicx}
\usepackage{pdfpages}
\usepackage{listings}
\lstset{
    breaklines=true,           % Enable word wrapping
    breakatwhitespace=false,   % Allows breaks even if there's no whitespace
}
\usepackage[a4paper, margin=0.125in, includefoot]{geometry} % Optional: Adjust margins as needed
\graphicspath{ {./Images/} }
% --- Preamble: Metadata ---
\title{}
\author{Justin K}
\date{\today} % Or use a specific date, e.g., \date{January 19, 2026}

% ---- Nacris ----
\newcommand{\As}{\emph{async}\xspace}
\newcommand{\Aw}{\emph{await}\xspace}
\newcommand{\Bbl}{\emph{Babel}\xspace}
\newcommand{\Bs}{\emph{Bootstrap}\xspace}
\newcommand{\Bz}{\emph{Blazor}\xspace}
\newcommand{\Cs}{\emph{C\#}\xspace}
\newcommand{\Cn}{\emph{continuation}\xspace}
\newcommand{\Ex}{\emph{Express}\xspace}
\newcommand{\Hk}{\emph{Hook}\xspace}
\newcommand{\Hks}{\emph{Hook}s\xspace}
\newcommand{\IA}{\emph{Interactive Auto}\xspace}
\newcommand{\IAM}{\emph{IAM}\xspace}
\newcommand{\IDP}{\emph{IDP}\xspace}
\newcommand{\IS}{\emph{Interactive Server}\xspace}
\newcommand{\IWA}{\emph{Interactive Web Assembly}\xspace}
\newcommand{\JS}{\emph{JavaScript}\xspace}
\newcommand{\JSX}{\emph{JSX}\xspace}
\newcommand{\Nd}{\emph{Node}\xspace}
\newcommand{\NJS}{\emph{NextJS}\xspace}
\newcommand{\OA}{\emph{OAuth2}\xspace}
\newcommand{\OIDC}{\emph{OpenID Connect}\xspace}
\newcommand{\Rt}{\emph{React}\xspace}
\newcommand{\PR}{\emph{Prop}\xspace}
\newcommand{\Prs}{\emph{Prop}s\xspace}
\newcommand{\RtD}{\emph{ReactDOM}\xspace}
\newcommand{\Rz}{\emph{Razor}\xspace}
\newcommand{\SPA}{\emph{SPA}\xspace}
\newcommand{\SPAs}{\emph{SPA}s\xspace}
\newcommand{\SSR}{\emph{SSR}\xspace}
\newcommand{\St}{\emph{State}\xspace}
\newcommand{\TS}{\emph{TypeScript}\xspace}
\newcommand{\Vt}{\emph{Vite}\xspace}
\newcommand{\WA}{\emph{WebAssembly}\xspace}
\newcommand{\UE}{\emph{useEffect}\xspace}

\begin{document}

% --- Title and Table of Contents ---
\maketitle
\tableofcontents
\begin{flushleft}


% ------------------------ Section ------------------------ 
\section{Initial Setup}

\begin{enumerate}
\item  dotnet new list 
\newline
Lists templates for starting new applications.
\item  You can view files in the project by using either the \emph{File Explorer} or the \emph{Solution Explorer}.
\item \emph{Program.cs} creates an instance of the \emph{app} object, which is an instance of the \emph{WebApplication} object that is the hosting application.
\item \emph{app} starts an HTTP server for your application so that you can start listening to HTTP requests.
\item The \emph{builder} object can be used to configure services for the application.
\item The \emph{Build()} method introduces the \emph{Kestrel Web Server}, reading \emph{appsettings.json}.
\item \verb|<Project Sdk="Microsoft.NET.Sdk.Web">|
\newline
The above from the \emph{.csproj} file imports all the libraries required for a web API application.
\item dotnet build
\newline
to build the application from your solution directory.
\item Use a \emph{record} for a \emph{Dto}, since it represents a contract between a client and a server as to how data will be transferred and this data should be immutable
\newline

\end{enumerate}

% ------------------------ Section ------------------------ 
\section{Entra Setup}

\begin{enumerate}
\item Log into Entra ID.
\item Select the tenant for the application. To begin with, set the organizational directory for the tenant as single-tenant. 
\item The tenant should have at least one user with the role "Application Developer".
\item The backend API will accept tokens that have been generate by only our tenant.
\item Every API that you register in Entra is going to need a \emph{scope} that defines the permissions required by a client that requires access to this API.
\item Scopes are defined by the \emph{Expose an API} section in Entra. every scope is prefixed with \emph{api://}.

\end{enumerate}

% ------------------------ Section ------------------------ 
\section{Securing the ASP.NET Core API in Entra}

\begin{enumerate}
\item Each endpoint that is restricted should have \emph{.RequireAuthorization()} appended.
\item The ASP.NET API back-end code needs to be able to validate the access tokens that will be attached to API requests that it receives. 
\newline 
This requires NuGet package \emph{Microsoft.AspNetCore.Authentication.JwtBearer}.
\item An \emph{Authentication Scheme} defines which handler to use to validate incoming access tokens.
\item Update the middleware in the ASP.NET Core application for authentication and authorization. 
\item To authorize the backend API to make requests, an access token will need to be attached to each request. To get this application token, a client application has to be registered with Entra.
\item The client application can be Postman, for testing purposes. The Redirect URI is \verb|https://oauth.pstmn.io/v1/callback|.
\newline
Register Postman as an App in Entra.
\item Entra has to be notified that Postman can talk to the ASP.NET Core back-end API.
\item Postman has to be granted API permissions in Entra to access the scope defined in the back-end API.
\end{enumerate}

% ------------------------ Section ------------------------ 
\section{Dependency Injection And Database Contexts}
\begin{enumerate}
\item DbContext is not thread-safe. Setting the lifetime of DbContext to \emph{Scoped} ensures that the connections do not span across HTTP requests.
\item DbContexts are an expensive resource. A scoped lifetime for DbContexts helps ensure that resources are used wisely.
\item \emph{Scoped} lifetime services are created once per HTTP request and reused within that HTTP request.
\item An environment variable can be set in PowerShell to hide the connection string during production, so that this string is not coded in the \emph{appsettings.json} file.
\begin{lstlisting}[breaklines=true]
$env:ConnectionStrings__DefaultConnection="DefaultConnection": "Data Source=localhost,1433;Persist Security Info=True;User ID=sa;Password=Mssql2025!;Pooling=False;MultipleActiveResultSets=False;Encrypt=True;TrustServerCertificate=True;Application Name=\"SQL Server Management Studio\";Command Timeout=0;Database=RESTAPI;"
\end{lstlisting}
The code in \emph{Program.cs} does not change. Setting the environment variable as shown above will add the connection string \emph{"ConnectionStrings:DefaultConnection"} to \emph{builder.Configuration}.

\end{enumerate}

\end{flushleft}

\end{document}
